\documentclass[11pt,a4paper]{article}
\usepackage[utf8]{inputenc}
\usepackage[margin=1in]{geometry}
\usepackage{amsmath,amssymb}
\usepackage{xcolor}
\usepackage{booktabs}
\usepackage{hyperref}
\usepackage{fancyhdr}
\usepackage{enumitem}

\definecolor{primary}{RGB}{31, 78, 121}
\definecolor{secondary}{RGB}{41, 128, 185}
\definecolor{darktext}{RGB}{44, 62, 80}

\pagestyle{fancy}
\fancyhf{}
\rhead{\textcolor{secondary}{\small HealthConnect Week 5 Project Summary}}
\lhead{\textcolor{secondary}{\small AnalystLab Africa}}
\rfoot{\textcolor{secondary}{\small Page \thepage}}
\lfoot{\textcolor{secondary}{\small Intern: Grai Rudolf (ML Track)}}

\begin{document}

\begin{center}
{\Large\bfseries\color{primary} ANALYSTLAB AFRICA EXPERIENCE LAB}\\[0.2cm]
{\large\bfseries\color{secondary} Official Week 5 Project Summary — Machine Learning Engineering Track}\\[0.4cm]
\rule{\linewidth}{1pt}
\end{center}

\vspace{0.2cm}

\begin{description}[font=\normalfont\bfseries\color{primary}]
    \item[Intern Name:] Grai Rudolf
    \item[Track:] Machine Learning Engineering / Data Science
    \item[Project Title:] HealthConnect Clinic Appointment Attendance Optimization
    \item[Date:] September 2026
\end{description}

\vspace{0.3cm}
\hrule
\vspace{0.5cm}

\section*{1. What I Planned to Accomplish in Week 5}
Building on the Week 4 problem definition and production system design, I planned to move from
planning into implementation by: constructing modular, reusable data-processing and preprocessing
pipelines; engineering behavioural features; training and benchmarking baseline classifiers on
recall-focused metrics; evaluating initial results; and testing the individual pipeline components.

\section*{2. What I Actually Completed}
\begin{enumerate}
    \item A modular \textbf{ML pipeline package} (\texttt{pipeline/}) with single-responsibility
          modules for configuration, data processing, feature engineering, model training and
          evaluation.
    \item A \textbf{data-processing workflow} with explicit leakage control (dropping
          \texttt{waiting\_time\_minutes}), NMAR/MCAR imputation, and a reusable sklearn
          \texttt{ColumnTransformer}.
    \item \textbf{Feature engineering} adding four behavioural features.
    \item An \textbf{initial model workflow} training and benchmarking \textbf{five} baselines on
          5-fold stratified cross-validation.
    \item An \textbf{evaluation framework} with metrics (accuracy, precision, recall, F1, ROC-AUC)
          and diagnostic visualisations.
    \item \textbf{Component testing} — all 6 smoke tests pass.
    \item Documents: \texttt{README.md}, \texttt{requirements.txt}, this report and a project
          summary (LaTeX/PDF), plus a fully executed Jupyter notebook.
\end{enumerate}

\section*{3. Key Findings / Development Outcomes}
\begin{itemize}
    \item The primary baseline (Logistic Regression) achieves held-out \textbf{recall $0.650$},
          \textbf{accuracy $0.617$}, and \textbf{ROC-AUC $0.664$}, comfortably above the random
          baseline and confirming pre-appointment signal.
    \item \texttt{booking\_lead\_days} dominates feature importance, followed by
          \texttt{distance\_to\_clinic\_km}, \texttt{age}, and appointment history — consistent with
          Week 4 findings.
    \item All five baselines are broadly comparable in recall ($0.56$–$0.63$), indicating a
          reliable starting point for tuning rather than a final model.
\end{itemize}

\section*{4. Major Challenges Encountered}
\begin{itemize}
    \item Designing preprocessing and imputation to preserve the meaning of NMAR missingness.
    \item Managing the dataset reduction (5,000 $\to$ 4,737 rows) after dropping cancellations.
    \item Handling the \texttt{booking\_lead\_days $= 0$} binning edge case to avoid NaN labels.
    \item Acknowledging that a stratified split does not fully control repeat-patient leakage
          (needs \texttt{GroupKFold}).
\end{itemize}

\section*{5. Important Decisions Made and Why}
\begin{itemize}
    \item \textbf{Selected logistic regression as the primary reporting baseline} for
          interpretability at this initial stage.
    \item \textbf{Used stratified split for the baseline} to establish performance quickly, with
          \texttt{GroupKFold} explicitly deferred to Week 6.
    \item \textbf{Used balanced class weighting} to favour no-show recall over precision, matching
          the business objective of catching high-risk patients.
\end{itemize}

\section*{6. Changes to My Week 4 Approach}
The Week 4 architecture and problem definition were followed unchanged. The main refinement was
choosing an interpretable linear baseline for the initial implementation while retaining tree-based
models for later feature-importance and hyperparameter work.

\section*{7. Cross-Track Collaboration}
Liaised with the Data Science / Data Analytics tracks to confirm the target definition (Attended vs
No-Show), the exclusion of cancellations, and the leakage boundaries. This ensured the engineered
features and pipeline align with the analytics outputs produced across the project team.

\section*{8. Remaining Work}
Notes: apply \texttt{GroupKFold} and a temporal split; tune hyperparameters and threshold for the
$F_2$ objective; containerise the inference service; add MLflow tracking and drift monitoring.

\section*{9. Proposed Focus for Week 6}
\begin{enumerate}
    \item Rigorous leakage-controlled validation (\texttt{GroupKFold} by \texttt{patient\_id}).
    \item Hyperparameter and threshold tuning for recall-driven objectives.
    \item Wrapping the best model in the containerised FastAPI inference service from the Week 4
          architecture.
    \item Introducing MLflow experiment tracking and Evidently AI drift monitoring.
\end{enumerate}

\end{document}
